\hspace{3mm}

\centerline{\bf \huge 7 класс}

\hspace{3mm} 

\problem{1} Эрдем поехал на Олимпиаду Физтех в Волгоград. По дороге из Элисты в Волгоград было 15 Макдональдсов и 14 KFC. В каждом фаст-фуд ресторане Эрдем покупал по чизбургеру и картошке фри, но из-за нехватки средств, после 9го Макдональдса Эрдем в оставшихся 12 ресторанах покупал только картошку. На обратном пути, ему скинули деньги на дорогу, и начиная с 6го Макдональдса он опять стал покупать по чизбургеру и картошке фри и ещё по коле, в оставшихся фаст-фуд ресторанах. Сколько денег потратил Эрдем на еду, если чизбургер стоит 50 рублей, кола стоит 45 рублей, а картофель фри стоит 30 рублей.

\begin{flushright}
    \scriptsize{$-$ Много мяса, нет салата — это завтрак для пирата}
\end{flushright}

\hspace{3mm}

\problem{2} Герман пришел на I ежегодную Математическую Олимпиаду ЭлМШ, на которой было предложено 5 задач. Герман очень торопился и мог потратить на Олимпиаду максимум час. Задачи, придуманные Расулом Владимировичем, Герман читает за 2 минуты, а решает за 35 минут. Задачи, придуманные Анджуром Аркадьевичем, он читает 2 минуты, а решает за 15 минут. Задачи, придуманные Очиром Владимировичем, он читает 10 минут и решает тоже за 10 минут. Очир Владимирович предложил две задачи, Анджур Аркадьевич одну, а Расул Владимирович две. Какое наибольшее количество баллов он успел набрать, если задачи Очира Владимировича он мог отличить, даже не читая, а отличить задачи Анджура Аркадьевича от задач Расула Владимировича может только прочитав их, при этом, читая, сначала ему попадаются задачи только Расула Владимировича. За задачи Очира Владимировича ученики получали по 4 балла, Анджура Аркадьевича - по 5 баллов, Расула Владимировича - по 8 баллов.

\begin{flushright}
    \scriptsize{$-$ Ты думаешь, что ты понял это,\\ хотя на самом деле всего лишь думаешь,\\ что ты понял то, о чем ты думаешь.}
\end{flushright}

\hspace{3mm}

\problem{3} Эвена и Эрдем играют в кретики-нолики, но не на классической доске $3\times 3$, а на бесконечной во все стороны причем Эвена ходит первая крестиками. Придумайте стратегию для Эвену как выиграть Эрдема причем не более чем за 6 ходов, если игра идёт до 4 знаков подряд (по вертикали или горизонтали).

\begin{flushright}
    \scriptsize{$-$ В этом мире нет справедливости.\\ Не удивительно, что он трещит по швам.}
\end{flushright}

\newpage

\hspace{3mm}

\problem{4} Божья коровка и муравей устроили забег на часах. Муравей бежал по минутной стрелке часов, а божья коровка по часовой стрелке. Когда впервые с начала гонки стрелки на часах пересеклись, муравей пробежал $\frac{2}{5}$ дистанции, а божья коровка только $\frac{1}{4}$. Когда божья коровка завершила забег, если муравей добежал до конца стрелки в 12:00? Забег начался одновременно, а божья коровка и муравей бежали с постоянной скоростью.

\begin{flushright}
    \scriptsize{$-$ Стрелки часов не повернуть вспять.\\ Но в наших силах двигать их вперед!}
\end{flushright}

\hspace{3mm}

\problem{5} В Республике Калмыкия любые два города связаны одним из трёх видов транспорта. Известно, что не существует города, обеспеченного сразу всеми тремя видами транспорта и так же известно, что не существует тройки городов, любые два из которых связаны одним и тем же видом транспорта. Найти наибольшее возможное количество городов в Калмыкии.

\begin{flushright}
    \scriptsize{$-$ Я не собираюсь жить тысячу лет.\\ Мне достаточно пережить сегодняшний день.}
\end{flushright}